\section{Opening and moving boundary pins}\label{sec:pins}

We transfer the plane reflection identity to two laws that bracket the
boundary problem.  Three separate operations are involved: opening small
gaps makes the pin sets into separated arcs; analytic translation moves
these arcs to one side of the insertion; and a local coupling closes the
gaps.  Throughout this section all geometric lengths are in physical
coordinates.  Every fixed arrangement is discretized before any of its
gaps or windows are shrunk.

Write $g_\delta=0$ on $A^+_\delta$ and $g_\delta=-1$ on
$A^-_\delta$.  A \emph{local increment function in a disk $K$} is a
measurable function of the increments on lattice edges in $K$, with a
fixed arbitrarily small enlargement of $K$ allowed for interpolation.
In particular, if $f\in C_c^\infty(K)$ and $\int f=0$, then
$H_\delta(f)$ is such a function: its interpolation coefficients have
sum zero, and paths joining their vertices can be chosen in the enlarged
disk.  We use
\begin{equation}\label{eq:pins-insertion}
 X_\delta=H_\delta(\Delta\phi),\qquad
 F_\delta=\prod_{j=1}^q F_{j,\delta},\qquad
 \supp\phi\Subset K_0,
\end{equation}
where $\phi\in C_c^\infty(K_0)$ is real and each $F_{j,\delta}$ is a local
increment function in a disk $K_j$.  We assume a uniform bound on
$\prod_j\|F_{j,\delta}\|_\infty$.  The disks $K_1,\ldots,K_q$
may overlap each other.  They are disjoint, with positive clearance,
from $K_0$.  The case $q=0$ means $F_\delta=1$.

\subsection{Windows and actual height brackets}

Fix a finite set $\mathcal L$ of horizontal levels.  Require that every
level meeting $\partial D$ is a regular value of the height coordinate
on $\partial D$, and that no level contains a marked point.  Its
intersections with $\partial D$ are finite: transversality makes them
isolated, and compactness then makes their number finite.  Around every
intersection choose a small open boundary arc, called a \emph{window}.
The closures of the windows are disjoint and avoid the marks.  Choose
them small enough that the portions of boundary between successive
windows contain nonempty exact-value arcs.  We denote this finite
window system by $\mathcal W$.

The discrete windows are intervals in the boundary parametrization of
$D_\delta$.  Uniform convergence of that parametrization specifies
them for all sufficiently small $\delta$; their endpoints are rounded
to partitions between consecutive boundary vertices.  Thus a vertex
is assigned to exactly one portion of the boundary.  Uniform
convergence, rather than the number of discrete crossings of a line,
is what determines these windows.

In the plane spin law impose the exact Gray-code pins of $g_\delta$
off $\mathcal W$.  In an upper bracket impose only $B=+$ on every
window, and in a lower bracket impose only $B=-$.  These finite pin
events have positive probability.  For example, a double cut around
the whole boundary reduces positivity to an admissible finite spin
extension; the corresponding bounded height extension supplies one.
Anchor the height at an exact-value vertex $v_*$ by
$H(v_*)=g_\delta(v_*)$.  Let $\mu^{+,\mathcal W}_\delta$ and
$\mu^{-,\mathcal W}_\delta$ be the resulting laws inside $D_\delta$.

\begin{lemma}[Boundary brackets]\label{lem:boundary-brackets}
The boundary traces in these laws satisfy the following actual integer
height constraints:
\begin{equation}\label{eq:bracket-intervals}
\begin{array}{c|cc}
 & \text{off the windows}&\text{on a window}\\ \hline
 \mu^{+,\mathcal W}_\delta & H=g_\delta&H\in\{0,1\}\\
 \mu^{-,\mathcal W}_\delta & H=g_\delta&H\in\{-2,-1\}.
\end{array}
\end{equation}
Conditional on the trace, the interior has the uniform height Gibbs
law with that trace.  Consequently there is a coupling with the true
height law $\mu^g_\delta$ such that
\begin{equation}\label{eq:bracket-order}
 H^-_\delta\le H^g_\delta\le H^+_\delta
 \quad\text{at every vertex, and hence after affine interpolation}.
\end{equation}
The bracket means lie between $-2$ and $1$.  The bounded-boundary
smearing estimates of \cref{prop:smearing-moments} apply to both
brackets, uniformly in the sizes and number of windows.
\end{lemma}

\begin{proof}
The inverse images of $B=+$ in $\Z$ are the disjoint adjacent pairs
$\{4t,4t+1\}$; those of $B=-$ are $\{4t-2,4t-1\}$.  Along a
connected chain with constant $B$, the height cannot move from one
such pair to another, since the distance between successive pairs is
three.  The nearest-neighbor height constraint therefore propagates
one pair along the entire chain.

For an upper window adjoining exact height $0$, this pair is
$\{0,1\}$.  If it adjoins exact height $-1$, its first vertex must
have height $0$, again selecting $\{0,1\}$.  For a lower window,
height $-1$ directly selects $\{-2,-1\}$, and an adjoining exact
height $0$ forces the first window vertex to be $-1$.  At a marked
transition the exact residues $0$ and $-1$ select the unique adjacent
branches.  Starting from $v_*$ and following the connected boundary
therefore gives \eqref{eq:bracket-intervals} everywhere, including
every later exact-value arc.  This proves an integer statement,
independent of any choice of representatives modulo four.

The spin-height bijection and the Gibbs property in
\cref{lem:gray} identify the conditional interior law.  No edge from
an interior vertex reaches past the pinned vertex boundary.  The
upper trace dominates $g_\delta$, and the lower trace is dominated
by it.  Height comparison, applied to each trace and then averaged,
gives the two stochastic orders and an ordered coupling.  Comparison
with constant boundary heights $-2$ and $1$, whose centered laws are
sign-symmetric, bounds the means.  The same comparison and
\cref{prop:smearing-moments} give the asserted moment bounds.
\end{proof}

\subsection{Reflection estimates with rare pins}

Analytic translation across a separating line is also used to propagate
correlation identities in the six-vertex argument of
\cite[Section~7.1, especially Lemma~7.4]{DCKLM2026}.  Our factors include
conditioning on rare obstacle events.  The first task is therefore to
bound the reflected norms after division by their probabilities; this
bound is what permits the same analytic-continuation principle here.

We first record precisely the normalization used when moving pins.
An obstacle event $P_i$ below consists of pins on a connected arc,
with the following description in a height lift:
\begin{equation}\label{eq:obstacle-intervals}
 H(v)\in I_i(v)+4t_i\quad(v\text{ on the arc}),\qquad t_i\in\Z.
\end{equation}
Here the $I_i(v)$ are fixed integer intervals, at least one of them
is a singleton, and the spin description consists of $B$ pins and
$W=+$ on a bounded number of connected subpaths.  These are the
obstacles of \cref{prop:pin-product,prop:obstacle-moments}.
The filled obstacle neighborhoods are pairwise disjoint, as specified
in \cref{prop:pin-product}.  Every unbounded insertion has a pin-free
filled test neighborhood, including its cut-search buffer, disjoint
from all of them, as in \cref{prop:obstacle-moments}.  These are fixed
positive-clearance conditions.  Bounded spectator disks may overlap one
another.  Constants may depend on this fixed geometry.

The disconnected pin arcs carry separate offsets $4t_i$.  We impose
their equality before closing gaps: \cref{prop:obstacle-moments} does
not give uniform constants as gaps shrink, whereas common offsets give
fixed integer interval constraints after anchoring.  Balanced endpoint
phases will Fourier-project the event that these offsets agree;
see \eqref{eq:fourier-offset-equality}.  To carry those phases through
reflections and translations, we use the auxiliary fields of
\cref{prop:reflection-positivity}.  For a real frequency $\theta$, set
\begin{equation}\label{eq:haar-phase}
 Z_\theta(v)=U_\theta
       \exp\{i\theta(H(v)-H(o))\},
\end{equation}
where $U_\theta$ is uniform on the unit circle, independent of the
spins; use an independent $U$ for each field being added.  The choice
of $o$ is immaterial because of Haar invariance, and the original spin
and increment law is unchanged.  \Cref{prop:reflection-positivity} gives
stationarity, reflection positivity, and the positive normal transfer
for the joint law.  Phase values in a half-plane are recovered from
its row phases and paths staying in that half-plane, so endpoint phases
have the required local support.  Frequencies may be arbitrary real
numbers; no periodicity of a lifted height is assumed.

Let $\vartheta$ be reflection across a lattice row.  On functions
supported in its closed positive half-plane, write
\begin{equation}\label{eq:pins-reflection-product}
 \langle A,B\rangle_{\vartheta}
     =\EE[\overline{\vartheta A}\,B],\qquad
 \|A\|_{\vartheta}^2=\langle A,A\rangle_{\vartheta}.
\end{equation}
We use the quotient and completion for this positive semidefinite
form.  Pure normal translation through
$\tau_\delta=\sqrt3\,\delta$ acts there as a positive self-adjoint
contraction $S_\delta$ by \cref{prop:reflection-positivity}.

\begin{lemma}[Normalized reflection bound]\label{lem:normalized-reflection}
Suppose the obstacle events $P_1,\ldots,P_m$ and the support of $A$
are strictly on one side of a reflection row.  The obstacles have
pairwise disjoint filled isolating neighborhoods, and every unbounded
insertion in $A$ has a pin-free filled test neighborhood disjoint from
them.  Choose these neighborhoods inside the same open half-plane.  Let
$p_i=\PP(P_i)$ and
\[
 V=A\prod_{i=1}^m\frac{\one_{P_i}}{p_i}.
\]
Assume $A$ is either uniformly bounded or a uniformly bounded
function times one smooth zero-total height insertion in such a
disk.  Bounded factors may include the phase fields
\eqref{eq:haar-phase}.  Then $\|V\|_{\vartheta}\le C$, uniformly
as $\delta\downarrow0$.  The assertion is uniform when the
macroscopic arrangements range over a compact family preserving
all the stated positive separations.
\end{lemma}

\begin{proof}
Let $P=\bigcap_iP_i$ and $J=\prod_i p_i$.  Directly from the
definition,
\begin{equation}\label{eq:squared-pin-normalization}
 \|V\|_{\vartheta}^{2}
 =\frac{\EE[\one_{P\cap\vartheta P}
                    \overline{\vartheta A}\,A]}{J^2}.
\end{equation}
There are $2m$ obstacles in this expression.  The strict separation
from the row separates the original obstacles from their reflected
copies; their individual probabilities are
$p_1,\ldots,p_m,p_1,\ldots,p_m$.  Consequently
\cref{prop:pin-product} gives
\begin{equation}\label{eq:reflected-pin-product}
 \PP(P\cap\vartheta P)\le C\prod_{i=1}^m p_i^2=CJ^2.
\end{equation}
For bounded $A$ this proves the claim.  Otherwise write $A=BY$,
with $B$ bounded and $Y$ the zero-total insertion.
\Cref{prop:obstacle-moments}, applied separately to $Y$ and
$\vartheta Y$ in the full reflected obstacle arrangement, gives
\[
 \EE[|Y\,\vartheta Y|\mid P\cap\vartheta P]
 \le\EE[|Y|^2\mid P\cap\vartheta P]^{1/2}
      \EE[|\vartheta Y|^2\mid P\cap\vartheta P]^{1/2}\le C.
\]
Combining this with \eqref{eq:reflected-pin-product} cancels the
\emph{squared} denominator in \eqref{eq:squared-pin-normalization}.
The estimates invoked have uniform constants on the indicated
compact families.  Unit-modulus phases do not alter them.
\end{proof}

\begin{lemma}[Analytic translation]\label{lem:analytic-translation}
Partition a finite arrangement of pin obstacles and local factors
into two groups strictly separated by projection on a lattice
reflection normal $n$.  The factors have the form allowed in
\cref{lem:normalized-reflection}, with at most one unbounded
zero-total insertion in the entire arrangement.  Move the high
group rigidly by $tn$, keeping the low group fixed, and normalize
the plane expectation by the product of the individual obstacle
probabilities.  Denote this quantity, at lattice normal step
translations, by $M_\delta(t)$.

If $M_\delta(t_\delta)\to0$ for every $t$ in a nonempty open
interval $I\subset(0,\infty)$ and lattice-rounded $t_\delta\to t$,
then $M_\delta(0)\to0$.  The same implication holds with the low
group moved in direction $-n$.  It also holds along sequences of
other fixed placements converging with all the required positive
clearances.
\end{lemma}

\begin{proof}
We prove the high-group version.  Choose $\eta>0$ so small that
moving this group by $-\eta n$ still leaves a fixed positive
projection gap.  Put $\eta_\delta=N_\delta \tau_\delta\to\eta$ with
$N_\delta\in\N$, and use the translated arrangement at
$-\eta_\delta$ as the initial arrangement.  Choose a lattice
reflection row in its gap, tending to a row strictly in the
limiting gap.  Reflect and conjugate the normalized low-group
factor to obtain a positive-side Hilbert vector $A_\delta$;
denote the normalized high-group factor at $-\eta_\delta$ by
$B_\delta$.  Individual obstacle probabilities are invariant under
these lattice translations and reflections.  By
\cref{lem:normalized-reflection},
\begin{equation}\label{eq:analytic-side-bound}
 \|A_\delta\|_{\vartheta},\ \|B_\delta\|_{\vartheta}\le C.
\end{equation}

For $\Rea z>-\eta_\delta$ define
\begin{equation}\label{eq:analytic-matrix-element}
 f_\delta(z)=
  \big\langle A_\delta,
  S_\delta^{(z+\eta_\delta)/\tau_\delta}B_\delta\big\rangle_{\vartheta}.
\end{equation}
If $S_\delta=\int_{[0,1]}\lambda\,dE_\delta(\lambda)$, the power
in this formula means
\[
 S_\delta^w=\int_{(0,1]}e^{w\log\lambda}\,dE_\delta(\lambda),
 \qquad \Rea w>0.
\]
Its value on the zero spectral subspace is zero.  It is holomorphic
in this open half-plane and has operator norm at most one.
Holomorphy follows, for example, by dominated differentiation on
compact subsets, using boundedness of
$\lambda^a|\log\lambda|^k$ for $a>0$.  Thus
\eqref{eq:analytic-side-bound} makes $f_\delta$ uniformly bounded.
For small $\delta$ all these functions are defined on the common
half-plane $\Rea z>-\eta/2$.  This open half-plane contains zero
in its interior.

At every lattice parameter $t=m\tau_\delta\ge0$ the exponent in
\eqref{eq:analytic-matrix-element} is the positive integer
$m+N_\delta$, so the transfer interpretation gives exactly
$f_\delta(t)=M_\delta(t)$.  There is no assertion here about the
spectral power at the initial boundary point $-\eta_\delta$.
Cauchy's estimate on compact subsets gives a uniform derivative
bound, so replacing a real parameter by its nearest lattice
parameter changes $f_\delta$ by $o(1)$.

Along any subsequence, the normal-family theorem supplies a further
subsequence converging locally uniformly to a holomorphic function
$f$.  The hypothesis makes $f$ zero on $I$.  The identity theorem
then makes it zero throughout the common half-plane, in particular
at zero.  Every subsequence has this property, proving the claim.
Choose the filled neighborhoods strictly inside the two separated
projection regions.  Rigid motion preserves their within-group
separations, while moving the groups apart preserves their projection
gap.  The initial negative backoff is smaller than that gap.  Reflected
copies likewise have disjoint filled neighborhoods by their positive
row margin.  These observations verify the geometric hypotheses of
the side norm bounds throughout the motion.  Convergent variations of
the other placements preserve those bounds and the common half-plane.  Replacing $n$ by $-n$
proves the low-group version.
\end{proof}

\subsection{Gaps and three directions of motion}

Choose two levels $\ell_-<\ell_+$ from $\mathcal L$ such that
$K_0$ lies strictly between them and neither line meets any
$K_j$, $j\ge1$.  All these separations have positive margins.
Inside each window based at an intersection with these two levels,
remove a still smaller boundary subarc from the pin set.  These are
the \emph{gaps}.  For now their sizes are fixed and positive.  They
are chosen to contain the boundary portions in a sufficiently thin
strip about their level, so that all remaining pins stay a positive
distance from the two lines.  Transversality and uniform boundary
convergence ensure this at small mesh, even if the approximating
boundary has many crossings inside a window.

The remaining connected pin arcs are the obstacle events
$P_1,\ldots,P_m$.  At these fixed positive gap sizes they approach
finitely many disjoint compact simple arcs.  Thin regular neighborhoods
with endcaps give pairwise disjoint filled isolating neighborhoods, as
in \cref{prop:pin-product}; they can avoid the fixed interior test disks.
Each arc has its buffered collar inside its own neighborhood.  Positive
separation from the chosen levels also lets the neighborhoods stay on
the same sides of those levels.  Each arc contains an exact-value subarc.  The branch argument in
\cref{lem:boundary-brackets}, now applied along an individual arc,
shows that it has exactly the interval description
\eqref{eq:obstacle-intervals}, with
\[
 I_i(v)=\{g_\delta(v)\}\ \text{off windows},\qquad
 I_i(v)=\{0,1\}\ \text{or}\ \{-2,-1\}\ \text{on windows}.
\]
Only the integer offset $t_i$ can differ between distinct arcs.
Take an exact-value representative $v_i$ on each arc and put
$g_i=g_\delta(v_i)$.

The geometric separation used below is the following.  Set
\begin{equation}\label{eq:three-normals}
 e=(0,1),\qquad u=(\sqrt3/2,1/2),\qquad
 u'=(-\sqrt3/2,1/2).
\end{equation}
These are normals to the three lattice reflection directions.
The obstacles and spectator disks lying in the middle band must
split into two families $L,R$, with a number $\kappa>0$ such that
\begin{equation}\label{eq:middle-projections}
 \sup_{x\in L}x\cdot u
    <\inf_{x\in K_0}x\cdot u-\kappa,
 \qquad
 \inf_{x\in R}x\cdot u
    >\sup_{x\in K_0}x\cdot u+\kappa.
\end{equation}
Here and below a supremum over a family means over the union of
its supports.  Empty families are omitted.  All the remaining
obstacles and spectator disks lie wholly above $\ell_+$ (family
$T$) or below $\ell_-$ (family $B$).  Overlapping spectator disks
belong to the same family; their factors will always move together.

\begin{proposition}[Identity with separated gaps]\label{prop:gapped-identity}
Fix windows, gaps and tests satisfying the preceding conditions.
Let $P=\bigcap_{i=1}^mP_i$ and
$J_\delta=\prod_{i=1}^m\PP(P_i)$.  For any product $\Psi_\delta$
of unit-modulus phase factors at the representatives $v_i$,
\begin{equation}\label{eq:gapped-phase-identity}
 \frac{\EE[X_\delta F_\delta\Psi_\delta\one_P]}{J_\delta}
       \longrightarrow0.
\end{equation}
The phases may have arbitrary fixed real frequencies and the
independent Haar lifts of \eqref{eq:haar-phase}.  The assertion is
valid for every convergent sequence of lattice-rounded placements
with these strict separations.

If $E$ is the event that all offsets in
\eqref{eq:obstacle-intervals} agree, then also
\begin{equation}\label{eq:gapped-equal-offset-identity}
 \frac{\EE[X_\delta F_\delta\one_{P\cap E}]}{J_\delta}
       \longrightarrow0.
\end{equation}
\end{proposition}

\begin{proof}
We first prove \eqref{eq:gapped-phase-identity}.  Include each
representative phase factor with its obstacle.  Move all supports
apart by the following sequence of one-sided translations.  Each
translation is across a strict projection gap and moves the entire
selected side rigidly.

First move $T$ by $a_+e$ and $B$ by $-a_-e$, successively.
These are legitimate moves across the original horizontal gaps.
Choose the positive distances sufficiently large that $T$ is above
$K_0$ in $u$-projection and $B$ is below it in that projection.
Together with \eqref{eq:middle-projections}, the two combined
families
\[
 Q_+=T\cup R,\qquad Q_-=B\cup L
\]
are now strictly above and below $K_0$ in $u$-projection.

Next move $Q_+$ by $bu$ and $Q_-$ by $-cu$, successively.
Choose $b$ large enough that all of $Q_+$ is strictly above $K_0$
in vertical coordinate, and strictly below $K_0$ in $u'$-projection.
This is possible because $u\cdot e=1/2$ and
$u\cdot u'=-1/2$.  Then choose $c$ large enough that all of
$Q_-$ lies strictly above $K_0\cup Q_+$ in $u'$-projection.
Each of these moves increases its separating $u$-gap.  No
separation within either combined family changes.

Finally move $Q_-$ by $du'$.  Its $u'$-projection is already
separated from every other support, so this is another legitimate
one-sided move.  Because $u'\cdot e=1/2$, a sufficiently large
$d$ puts all of $Q_-$ strictly above $K_0$ vertically.  Thus every
factor other than $X_\delta$ is now strictly above $K_0$.
All choices can be made in open ranges with positive margins;
every finite choice leaves a fixed finite arrangement.  The
successive strict projection gaps also keep different moving
families disjoint, while distances inside each moving family are
preserved.  In particular, all individual arc collars remain
available.  \Cref{fig:pin-moving} illustrates the three directions.

At the final placement put a horizontal reflection row between
$K_0$ and all the other supports.  Reflection Cauchy--Schwarz and
\cref{lem:normalized-reflection} bound the normalized expectation
by
\[
 C\|X_\delta\|_{\vartheta}.
\]
The norm tends to zero by \cref{cor:reflection-null}.  That
corollary remains valid in the phase-augmented Hilbert space,
since $X_\delta$ itself does not involve phases and its reflection
inner product is unchanged by their addition.

For every fixed admissible $a_+,a_-,b,c$, this proves vanishing
for $d$ in a sufficiently large open interval.  Apply
\cref{lem:analytic-translation} to remove the last move.  The same
argument then applies for every sufficiently large $c$ with
$a_+,a_-,b$ fixed, so it removes the $-cu$ move.  Remove $bu$,
$-a_-e$ and $a_+e$ in that order, each time using an open interval
of allowed large values.  This is successive one-variable
analytic continuation; no simultaneous motion through surrounding
pins is involved.  It proves \eqref{eq:gapped-phase-identity} at
the original placement.

To impose $E$, note that on $P$
\[
 D_i=\frac{H(v_i)-g_i-H(v_1)+g_1}{4}=t_i-t_1\in\Z
 \quad(2\le i\le m).
\]
For $s=(s_2,\ldots,s_m)\in[0,2\pi]^{m-1}$,
\begin{equation}\label{eq:balanced-offset-phases}
 e^{i\sum_{i=2}^m s_iD_i}
 =\prod_{i=2}^m
   e^{-is_i(g_i-g_1)/4}
   Z_{s_i/4}^{(i)}(v_i)
   \overline{Z_{s_i/4}^{(i)}(v_1)}.
\end{equation}
Use an independent phase field for each $i$ on the right.  Every
Haar factor cancels with its conjugate.  Hence this is an identity
of functions of the original increments, including for unbounded
relative lifts.  Fourier orthogonality gives exactly
\begin{equation}\label{eq:fourier-offset-equality}
 \one_E=(2\pi)^{-(m-1)}
   \int_{[0,2\pi]^{m-1}}e^{i\sum_{i=2}^m s_iD_i}\,ds
 \quad\text{on }P.
\end{equation}
For each fixed $s$, the normalized expectation tends to zero by
\eqref{eq:gapped-phase-identity}.  Its absolute value is bounded
uniformly in $s$ and $\delta$: indeed
\[
 \frac{\EE[|X_\delta F_\delta|\one_P]}{J_\delta}
 \le \frac{\PP(P)}{J_\delta}\,
       \|F_\delta\|_\infty\,
       \EE[|X_\delta|\mid P]\le C
\]
by \cref{prop:pin-product,prop:obstacle-moments} at this fixed
gap geometry.  Dominated convergence in
\eqref{eq:fourier-offset-equality} proves
\eqref{eq:gapped-equal-offset-identity}.
\end{proof}

\begin{figure}[t]
 \centering
 \begin{tikzpicture}[x=1cm,y=1cm,>=Stealth,
  group/.style={draw=blue!55!black,fill=blue!8,rounded corners=3pt,
                minimum width=8mm,minimum height=6mm,font=\small},
  ins/.style={circle,draw=black,fill=white,inner sep=2pt,font=\small},
  move/.style={->,very thick,draw=blue!65!black},
  guide/.style={dashed,gray!55}]
\begin{scope}[shift={(0,0)}]
 \draw[guide] (-1.25,.55)--(1.25,.55);
 \draw[guide] (-1.25,-.55)--(1.25,-.55);
 \node[ins] at (0,0) {$X$};
 \node[group] at (-1,0) {$L$};
 \node[group] at (1,0) {$R$};
 \node[group] at (0,.95) {$T$};
 \node[group] at (0,-.95) {$B$};
 \draw[move] (0,1.32)--(0,1.88);
 \draw[move] (0,-1.32)--(0,-1.88);
 \node[font=\small,align=center] at (0,-2.4)
    {1. Vertical separation\\$+e$ and $-e$};
\end{scope}
\begin{scope}[shift={(4.5,0)}]
 \draw[guide] (-.7,1.22)--(.7,-1.22);
 \node[ins] at (0,0) {$X$};
 \node[group] at (.92,.56) {$T,R$};
 \node[group] at (-.92,-.56) {$B,L$};
 \draw[move] (1.35,.82)--(1.95,1.17)
    node[right,font=\small] {$u$};
 \draw[move] (-1.35,-.82)--(-1.95,-1.17);
 \node[font=\small,align=center] at (0,-2.4)
    {2. Oblique separation\\$+u$ and $-u$};
\end{scope}
\begin{scope}[shift={(10,0)}]
 \node[ins] at (0,0) {$X$};
 \node[group] at (1.04,1.1) {$T,R$};
 \node[group] at (-.65,-.7) {$B,L$};
 \draw[move] (-1.06,-.35)--(-2.792,.65)
    node[above,font=\small] {$u'$};
 \draw[guide] (-2.15,.48)--(1.55,.48);
 \node[font=\small,align=center] at (0,-2.4)
    {3. Raise the remaining group\\along $+u'$};
\end{scope}
\end{tikzpicture}


 \caption{The three normal directions used to move the obstacles.
 Each shaded group denotes all of its arc obstacles and local test
 factors.  The insertion $X$ stays fixed.  Vertical moves first
 separate $T$ and $B$; the $u$ moves separate $T\cup R$ and
 $B\cup L$; for a sufficiently large final displacement, the
 $u'$ move places the latter above $X$.
 The relation $u\cdot u'=-1/2$ supplies the projection gap for
 the last move.  Distances are chosen successively; the drawing is
 schematic.}\label{fig:pin-moving}
\end{figure}

\subsection{Equal offsets and filling the gaps}

Keep $\mathcal W$ fixed.  Around each selected intersection choose
a small physical neighborhood contained in its single-color
window, disjoint from all test disks and from all other such
neighborhoods.  In particular there are no $W$ pins in these
neighborhoods.  Let $r>0$ be a sufficiently small fixed size for
them.  Shrink the gaps so that each lies within a ball of radius
$a\ll r$ about its intersection.  This is possible by
transversality.  Constants in what follows can depend on
$\mathcal W$ and $r$, but not on $a$ or on sufficiently small
$\delta$.  We write $o_\delta(1)$ with $a$ fixed.

Let $\mu_{\delta,a}$ be the plane law conditioned on the gapped
pins.  Let $\mu^{=}_{\delta,a}=\mu_{\delta,a}(\,\cdot\mid E)$,
whenever $E$ has positive probability.  Finally let
$\mu_{\delta,0}$ denote the filled bracket.  All these notations
refer to one chosen sign of the bracket.

\begin{lemma}[Gap filling]\label{lem:gap-filling}
There are $\beta>0$ and a function $\epsilon(a)\to0$ as
$a\downarrow0$ such that
\begin{equation}\label{eq:offset-high-probability}
 \limsup_{\delta\downarrow0}\mu_{\delta,a}(E^c)
        \le\epsilon(a).
\end{equation}
If $\mathscr I_\delta$ is the joint increment information in the
fixed test disks, then
\begin{equation}\label{eq:gap-tv}
 \limsup_{\delta\downarrow0}
 \big\|\mu^{=}_{\delta,a}|_{\mathscr I_\delta}
          -\mu_{\delta,0}|_{\mathscr I_\delta}\big\|_{\rm TV}
       \le\epsilon(a).
\end{equation}
For every finite $p$, and for all sufficiently small $a$, the
variables $X_\delta$ have uniformly bounded $p$th moments under
$\mu^{=}_{\delta,a}$ and under $\mu_{\delta,0}$, in the sense
of a bound on $\limsup_{\delta\downarrow0}$ independent of $a$.
No such uniform assertion is needed for the law conditioned on
$E^c$.
\end{lemma}

\begin{proof}
First replace each requirement $W=+$ by the requirement that $W$
be constant on that subpath, and call the resulting edge equalities
$E_0$.  If there are $M$ prescribed subpaths, then, conditional on
$B$ and $E_0$, requiring all their component labels to be plus
has probability between $2^{-M}$ and one.  The number $M$ is
fixed by the original windows and does not grow as gaps shrink.

Consider one gap.  In the annulus from scale $Ca$ to scale $r$,
the only pins are $B$ pins of the desired window sign.  There are
no $E_0$ edges there.  The favorable-pin version of
\cref{lem:buffer-circuits}, used in a logarithmic collection of
separated annular slots, therefore gives a double circuit of
that sign surrounding the gap except with probability
\[
 C(a/r)^\beta+o_\delta(1).
\]
This estimate holds before any offset equality has been imposed.
The two adjoining pinned boundary arms leave the neighborhood,
so a surrounding dual circuit crosses each of them.  At a crossing
both incident primal sites have the window sign.  Consecutive
edges of a double circuit also share a connected chain of sites
of that sign.  Thus the circuit joins the two pinned arms through
a connected constant-$B$ chain.  The adjacent-pair argument of
\cref{lem:boundary-brackets} forces the same lifted pair on both
arms.  Their offsets are therefore equal.

The gap adjacency graph of the remaining boundary arcs is a
cycle, hence connected.  Circuits around all the gaps force every
$t_i$ to agree.  A union bound over the fixed number of gaps,
followed by restoration of the $W=+$ labels at cost at most
$2^M$, proves \eqref{eq:offset-high-probability}, with for
example a constant multiple of $(a/r)^\beta$ as the error.
In particular $\mu_{\delta,a}(E)$ is bounded away from zero for
small $a$ and then small $\delta$.

For the coupling assertion first work with $E_0$ and fill only
one gap.  All the new $B$ pins have the same sign.  The $B$
marginals before and after adding them are ordered, by
\cref{lem:spin-fkg}.  Use the stopping coupling of \cref{sec:cuts}, based on
\cref{lem:double-cut}, exploring outwards from the gap and
coupling the queried values in this order.  In the more adverse
law a double circuit of the added sign occurs before the search
leaves the window neighborhood with failure probability at most
$C(a/r)^\beta+o_\delta(1)$, by the same favorable-pin estimate.
On success it is a common double circuit in the two laws.  All
changed constraints are on its gap side; no $E_0$ edge crosses it.
The conditional law on its other side is therefore identical.
Sample that side identically, including its $W$ components.

Do this successively at the finitely many gaps.  Searches stay in
disjoint window neighborhoods.  The resulting coupling matches
the joint increments in every test disk and all the $W$ labels
that must subsequently be checked, except with probability
tending to zero with $a$.  This last point uses that the label
sites are outside the single-color window neighborhoods.  If two
laws can be coupled with error $\varepsilon$ on variables
including an event $L$, and both give $L$ probability at least
$c>0$, their restrictions conditioned on $L$ differ in total
variation by at most $2\varepsilon/c$.  Indeed, for any event
$A$ the numerators $\PP(A\cap L)$ differ by at most
$\varepsilon$, as do the two denominators, and division gives
the stated bound.  Applying this with $c=2^{-M}$ restores all
the required plus labels.  It proves the claimed approximation
between the unconditioned-offset gapped law and the filled law.
Conditioning the former additionally on $E$ changes any marginal
by at most $\mu_{\delta,a}(E^c)$.  This proves
\eqref{eq:gap-tv}, after enlarging $\epsilon(a)$.

It remains to check the uniform moments under equal offsets.
Fix an annular buffer around $K_0$, with fixed inner and outer
radii and positive clearance from $K_0$.  Choose its filled
interior compactly contained in $D$ and disjoint from all gap
neighborhoods, so the buffer and its interior contain no pins.
Conditional just on the pins, a
double $B$ cut in this buffer has probability at least a fixed
$c_0>0$.  This follows from the local cut estimate and the
bounded cost of the plus labels.  By
\eqref{eq:offset-high-probability}, its probability under
additional conditioning on $E$ is at least $c_0/2$ for all
sufficiently small $a$ and then small mesh.  Moreover $E$ is
measurable outside this cut: relative heights between the
representatives can be evaluated on paths in the connected
exterior of a disk containing the buffer.  Such paths avoid the
unrevealed interior of the cut.  The usual outside-in cut
exploration therefore leaves a Gibbs interior with its
bounded-range trace, even after $E$ is imposed.

All offsets are now common.  Anchoring one representative makes
the interval restrictions fixed integer intervals, so rounded
midpoints require no parity classes.  Together with the fixed
buffer, these observations verify the hypotheses of
\cref{cor:equal-offset-moments}; it
gives the asserted bounds uniformly as the gaps close.  The
filled law has bounded boundary data by
\cref{lem:boundary-brackets}, and
\cref{prop:smearing-moments} gives its bounds.  This argument
does not estimate the exceptional unequal-offset law.
\end{proof}

\begin{proposition}[Bracket identity]\label{prop:bracket-identity}
Fix a window system $\mathcal W$ for a finite grid
$\mathcal L$, and either bracket from
\cref{lem:boundary-brackets}.  Let $K_0,\ldots,K_q\Subset D$
be disks, with $K_0$ at positive distance from every $K_j$,
$j\ge1$, and let $X_\delta,F_\delta$ be as in
\eqref{eq:pins-insertion}.  Suppose there are two levels
$\ell_-<\ell_+$ in $\mathcal L$ such that:
\begin{enumerate}[label=\textnormal{(\roman*)}]
\item $K_0$ is strictly between the two levels, and both levels
      avoid every spectator disk $K_j$ with positive margin;
\item after opening sufficiently small gaps inside their
      corresponding windows, every pin arc and spectator disk
      in the middle band belongs to one of the two strictly
      separated $u$-projection families in
      \eqref{eq:middle-projections}.
\end{enumerate}
Then
\begin{equation}\label{eq:bracket-identity}
 \lim_{\delta\downarrow0}
   \EE_{\mu^{\pm,\mathcal W}_\delta}
       [H_\delta(\Delta\phi)F_\delta]=0.
\end{equation}
In particular it applies to products of bounded continuous
functions of finitely many smooth zero-total tests in each
$K_j$.  There is no disjointness requirement between spectator
disks.  For this latter class the identity passes to any
subsequential distributional limit of the bracket laws.

For any fixed bound $q$ on the number of spectator disks these
geometric hypotheses can be arranged locally as follows.  Fix
$x\in D$ and $\rho>0$ with $\overline{B(x,\rho)}\subset D$.
Take $K_0$ sufficiently small about $x$, and take at most $q$
spectator disks of sufficiently small radii, all outside
$B(x,\rho)$.  From a sufficiently dense grid of regular levels
one can choose the two levels within a fixed small fraction of
$\rho$ above and below $x$.  The choices and smallness bounds
depend on $q$ and $\rho$, not on the number of terms in a later
linear combination of such products.
\end{proposition}

\begin{proof}
At fixed gaps, \cref{prop:pin-product} gives
$\PP(P)/J_\delta\ge c_a>0$.  By
\cref{lem:gap-filling}, for all sufficiently small fixed $a$,
$\PP(E\mid P)\ge1/2$ at small mesh.  Division of
\eqref{eq:gapped-equal-offset-identity} by
$\PP(P\cap E)/J_\delta$ therefore gives
\begin{equation}\label{eq:conditional-gapped-identity}
 \lim_{\delta\downarrow0}
    \EE_{\mu^{=}_{\delta,a}}[X_\delta F_\delta]=0.
\end{equation}
The positive number $c_a$ may tend to zero as $a\downarrow0$;
it is used only with $a$ fixed.

Let $X^{(T)}=\max(-T,\min(X,T))$ and put
$M=\sup_\delta\|F_\delta\|_\infty$.  The total variation bound
in \cref{lem:gap-filling} applies to $X^{(T)}F_\delta$, which
is a bounded function of the joint test increments.  The same
lemma, at any $p>1$, bounds the truncation errors under the
equal-offset and filled laws by $C_pM T^{1-p}$.  Consequently
\[
 \limsup_{\delta\downarrow0}
 \left|\EE_{\mu_{\delta,0}}[X_\delta F_\delta]\right|
 \le 2MT\epsilon(a)+C_pM T^{1-p}.
\]
First send $a\downarrow0$ and then $T\to\infty$.  This proves
\eqref{eq:bracket-identity}.  For bounded continuous functions
of smooth tests, distributional convergence and the same
uniform integrability pass the expectation to a subsequential
limit.

For completeness, here is a quantitative choice of the local
geometry in the last assertion.  Translate $x$ to the origin.
Take $K_0\subset B(0,\rho/1000)$ and seek levels in
$(\rho/40,\rho/20)$ and $(-\rho/20,-\rho/40)$.
Choose spectator radii so small that the total length of their
vertical projections is less than half the length of either
candidate interval; a bound $\rho/(1000(q+1))$ on each radius
suffices.  There is then a nonempty open subinterval of each
candidate interval avoiding all spectator disks.  Regular values
are dense, so a sufficiently dense regular grid supplies both
levels, with positive margins.

Every boundary point in the resulting middle band is outside
$B(0,\rho)$.  The same is true of every middle spectator disk
by assumption.  Since $|y|<\rho/20$ there, each such point has
$|x|>\sqrt{1-1/400}\,\rho$.  Its $u$-projection
$(\sqrt3/2)x+y/2$ consequently has magnitude greater than
$\rho/2$, with the sign of $x$.  A connected middle arc or
disk cannot change that sign.  Thus
\eqref{eq:middle-projections} holds with ample margin, whereas
all projections of $K_0$ have magnitude at most $\rho/1000$.
Small gaps and lattice approximation errors preserve the
strict inequalities.  This proves the geometric assertion.
\end{proof}

The order of limits is now explicit.  The identity for a fixed
filled window system is proved by first sending the mesh to
zero at fixed positive gaps and then shrinking those gaps.
In the boundary comparison below, one subsequently shrinks
the windows for each fixed finite grid, and only then refines
the grid.  The bracket order alone does not identify its law
with the true boundary law; the next section supplies the
boundary information needed for that step.
